

\centerline{\large\bf Base-change without Flatness}

\medskip

The following lemma on base-change does not need any flatness hypothesis.
The price paid is that the integer $r_0$ may depend on $\phi$.


\begin{lemma}\label{nhq-base-change-without-flatness}\dcref{3.1}\source{nitsure-hilbert-quot:page-0013}
Let $\phi : T\to S$ be a morphism of noetherian schemes,
let $\F$ a coherent sheaf on $\PP^n_S$, and
let $\F_T$ denote the pull-back of
$\F$ under the induced morphism $\PP^n_T\to \PP^n_S$. Let
$\pi_S : \PP^n_S \to S$ and $\pi_T : \PP^n_T \to T$
denote the projections. Then there exists an integer
$r_0$ such that the base-change homomorphism
$$\phi^* {\pi_S}_*\, \F(r) \to {\pi_T}_* \, \F_T(r)$$
is an isomorphism for all $r\ge r_0$.
\end{lemma}



\begin{proof} As base-change holds for open embeddings, using a finite affine open
cover $U_i$ of $S$ and a finite affine open cover $V_{i,j}$ of each
$\phi^{-1}(U_i)$ (which is possible by noetherian hypothesis),
it is enough to consider the case where $S$ and
$T$ are affine.

Note that for all integers $i$, the base-change homomorphism
$$\phi^*{\pi_S}_*{\mathcal O}_{\PP^n_S}(i) \to {\pi_T}_* {\mathcal O}_{\PP^n_T}(i)$$
is an isomorphism. Moreover, if $a$ and $b$ are any integers
and if $f : {\mathcal O}_{\PP^n_S}(a) \to {\mathcal O}_{\PP^n_S}(b)$ is any homomorphism
and $f_T : {\mathcal O}_{\PP^n_T}(a) \to {\mathcal O}_{\PP^n_T}(b)$ denotes its pull-back
to $\PP^n_T$, then for all $i$ we have the following commutative diagram
where the vertical maps are base-change isomorphisms.
$$\begin{array}{ccc}
\phi^*{\pi_S}_*{\mathcal O}_{\PP^n_S}(a+i)
& \stackrel{\phi^*{\pi_S}_*f(i)}{\to} &
\phi^*{\pi_S}_*{\mathcal O}_{\PP^n_S}(b+i) \\
\dna & & \dna \\
{\pi_T}_* {\mathcal O}_{\PP^n_T}(a+i)
& \stackrel{{\pi_T}_*f_T(i)}{\to} &
{\pi_T}_* {\mathcal O}_{\PP^n_T}(b+i)
\end{array}
$$


As $S$ is noetherian and affine, there exists an exact sequence
$$\oplus^p {\mathcal O}_{\PP^n_S}(a)\, \stackrel{u}{\to}  \,
\oplus^q {\mathcal O}_{\PP^n_S}(b) \,\stackrel{v}{\to} \,\F \,\to \,0$$
for some integers $a$, $b$, $p\ge 0$, $q\ge 0$.
Its pull-back to $\PP^n_T$ is an exact sequence
$$\oplus^p {\mathcal O}_{\PP^n_T}(a)\, \stackrel{u_T}{\to}  \,
\oplus^q {\mathcal O}_{\PP^n_S}(b) \,\stackrel{v_T}{\to} \,\F_T \,\to \,0$$
Let $\G = \ker(v)$ and let $\H = \ker(v_T)$.
For any integer $r$, we get exact sequences
$${\pi_S}_* \oplus^p {\mathcal O}_{\PP^n_S}(a+r) \,\to \,
  {\pi_S}_* \oplus^q {\mathcal O}_{\PP^n_S}(b+r) \,\to \,
  {\pi_S}_* \F(r)                       \,\to \,
  R^1{\pi_S}_* \G(r)$$
and
$${\pi_T}_* \oplus^p {\mathcal O}_{\PP^n_T}(a+r)   \,\to \,
  {\pi_T}_* \oplus^q {\mathcal O}_{\PP^n_T}(b+r) \,\to \,
  {\pi_T}_* \F_T(r)                    \,\to \,
  R^1{\pi_T}_* \H(r)$$
There exists an integer $r_0$ such that
$R^1{\pi_S}_*\G(r) =0$ and
$R^1{\pi_T}_*\H(r)=0$ for all $r\ge r_0$.
Hence for all $r\ge r_0$, we have exact sequences
$${\pi_S}_* \oplus^p {\mathcal O}_{\PP^n_S}(a+r) \,\stackrel{{\pi_S}_*u(r)}{\to}\,
\oplus^q {\mathcal O}_{\PP^n_S}(b+r) \,\stackrel{{\pi_S}_*v(r)}{\to} \,
{\pi_S}_* \F(r) \,\to \, 0$$
and
$${\pi_T}_* \oplus^p {\mathcal O}_{\PP^n_T}(a+r) \,\stackrel{{\pi_T}_*u_T(r)}{\to}\,
{\pi_T}_*   \oplus^q {\mathcal O}_{\PP^n_T}(b+r) \,\stackrel{{\pi_T}_*v_T(r)}{\to} \,
{\pi_T}_* \F_T(r) \,\to\, 0$$


Pulling back the second-last exact sequence under $\phi: T\to S$,
we get the commutative diagram with exact rows
$$\begin{array}{ccccc}
\phi^*{\pi_S}_* \oplus^p {\mathcal O}_{\PP^n_S}(a+r) &
\stackrel{\phi^*{\pi_S}_*u}{\to} &
\phi^* \oplus^q {\mathcal O}_{\PP^n_S}(b+r) &
\stackrel{\phi^*{\pi_S}_*v}{\to} &
\phi^*{\pi_S}_* \F(r) \to  0 \\
\dna &   & \dna & & \dna ~~~~~~~~ \\
{\pi_T}_* \oplus^p {\mathcal O}_{\PP^n_T}(a+r) &
\stackrel{{\pi_T}_*u_T(r)}{\to} &
{\pi_T}_*   \oplus^q {\mathcal O}_{\PP^n_T}(b+r) &
\stackrel{{\pi_T}_*v_T(r)}{\to} &
{\pi_T}_* \F_T(r) ~~ \to  0
\end{array}
$$
in which the first row is exact by the right-exactness of tensor product.
The vertical maps are base-change homomorphisms, the first two
of which are isomorphisms for all $r$.
Therefore by the five lemma,
$\phi^*{\pi_S}_* \F(r) \to {\pi_T}_* \F_T(r)$ is an isomorphism
for all $r\ge r_0$.
\end{proof}


\medskip



{\footnotesize
The following elementary proof of the above result is taken from Mumford [M]:
Let $M$ be the graded ${\mathcal O}_S$-module
$\oplus_{m\in \Z}\, {\pi_S}_* \F(m)$, so that $\F = M^{\sim}$.
Let $\phi^*M$ be the graded ${\mathcal O}_T$-module which is the pull-back of
$M$. Then we have $\F_T = (\phi^*M)^{\sim}$.
On the other hand, let $N = \oplus_{m\in \Z}\, {\pi_T}_* \F_T(m)$,
so that we have $\F_T = N^{\sim}$. Therefore, in the category of
graded ${\mathcal O}_T[x_0,\ldots,x_n]$-modules, we get an induced equivalence
between $\phi^*M$ and $N$, which means the natural homomorphisms of
graded pieces $(\phi^*M)_m \to N_m$ are isomorphisms for all $m\gg 0$.
$\square$
}%end of footnotesize





\bigskip


\centerline{\large\bf Flatness of $\F$ from Local Freeness of $\pi_*\F(r)$}


\medskip

\begin{lemma}\label{nhq-graded-module-flat}\dcref{3.2}\source{nitsure-hilbert-quot:page-0015}
Let $S$ be a noetherian scheme and let $\F$ be a coherent sheaf on
$\PP^n_S$. Suppose that there exists some integer $N$ such that
for all $r \ge N$ the direct image $\pi_*\F(r)$ is locally free.
Then $\F$ is flat over $S$.
\end{lemma}

\begin{proof} Consider the graded module $M = \oplus_{r\ge N} M_r$ over
${\mathcal O}_S$, where $M_r =  \pi_*\F(r)$. The sheaf $\F$ is isomorphic
to the sheaf $M^{\sim}$ on
$\PP^n_S = {\bf Proj}_S\,{\mathcal O}_S[x_0,\ldots,x_n]$ made from the
graded sheaf $M$ of ${\mathcal O}_S$-modules.
As each $M_r$ is flat over ${\mathcal O}_S$, so is $M$. Therefore
for any $x_i$ the localisation $M_{x_i}$ is flat over ${\mathcal O}_S$.
There is a grading on $M_{x_i}$, indexed by $\Z$, defined
by putting $\deg(v_p/x_i^q) = p-q$ for $v_p \in M_p$ (this is well-defined).
Hence the component $(M_{x_i})_0$ of degree zero,
being a direct summand of $M_{x_i}$, is again
flat over ${\mathcal O}_S$. But by definition of $M^{\sim}$, this is just
$\Gamma(U_i,\F)$, where $U_i = {\bf Spec}_S\,{\mathcal O}_S[x_0/x_i,\ldots,x_n/x_i]
\subset \PP^n_S$. As the $U_i$ form an open cover of $\PP^n_S$,
it follows that $\F$ is flat over ${\mathcal O}_S$. \end{proof}


\medskip

\centerline{\large\bf Grothendieck Complex for Semi-Continuity}

\medskip

The following is a very important basic result of Grothendieck,
and the complex $K^{\cdot}$ occurring in it is called the Grothendieck
complex.

%%%For a good exposition, see Mumford's {\sl Abelian Varieties}.

\begin{theorem}\label{nhq-grothendieck-complex}\dcref{3.3}\source{nitsure-hilbert-quot:page-0015}
Let $\pi: X\to S$ be a proper morphism of noetherian
schemes where $S=\spec A$ is affine,
and let $\F$ be a coherent ${\mathcal O}_X$-module
which is flat over ${\mathcal O}_S$. Then
there exists a finite complex
$$0\to K^0 \to K^1 \to \ldots \to K^n\to 0$$
of finitely generated projective $A$-modules, together with
a functorial $A$-linear isomorphism
$$H^p(X, \F \otimes_A M)
\stackrel{\sim}{\to} H^p(K^{\cdot}\otimes_A M)$$
on the category of all $A$-modules $M$.
\end{theorem}

\begin{proof}
\notready The source states this foundational theorem without giving its proof.
\end{proof}


The above theorem is the foundation for all results about direct images and
base-change for flat families of sheaves, such as Theorem
\ref{nhq-base-change-flat-sheaves}.

As another consequence of the above theorem, we have the following.


\begin{theorem}\label{nhq-cokernel-dual}\dcref{3.4}\source{nitsure-hilbert-quot:page-0016}\uses{nhq-grothendieck-complex} {\rm ([EGA] III 7.7.6) }
Let $S$ be a noetherian scheme and $\pi : X \to S$ a proper morphism.
Let $\F$ be a coherent sheaf on $X$ which is flat over $S$.
Then there exists a coherent sheaf $\Qq$ on $S$ together with
a functorial ${\mathcal O}_S$-linear isomorphism
$$\theta_{\G} : \pi_*(\F \otimes_{{\mathcal O}_X} \pi^*\G) \to \hom_{{\mathcal O}_S}(\Qq,\G)$$
on the category of all quasi-coherent sheaves $\G$ on $S$.
By its universal property,
the pair $(\Qq, \theta)$ is unique up to a unique isomorphism.
\end{theorem}


\begin{proof} If $S =\spec A$, then
we can take $\Qq$ to be the coherent sheaf associated to the $A$-module
$Q$ which is the cokernel of the
transpose $\pa^{\vee} : (K^1)^{\vee} \to (K^0)^{\vee}$ where
$\pa: K^0 \to K^1$ is the differential of any chosen Grothendieck
complex of $A$-modules
$0\to K^0 \to K^1 \to \ldots \to K^n\to 0$ for the
sheaf $\F$, whose existence is given by Theorem \ref{nhq-grothendieck-complex}.
For any $A$-module $M$,
the right-exact sequence
$(K^1)^{\vee} \to (K^0)^{\vee} \to Q\to 0$ with $M$
gives on applying $Hom_A(-,M)$ a left-exact sequence
$$0 \to Hom_A(Q, M) \to K^0\otimes_AM \to K^1\otimes_AM$$
Therefore by Theorem \ref{nhq-grothendieck-complex}, we have
an isomorphism
$$\theta^A_M : H^0(X_A, \F_A \otimes_A M)  \to  Hom_A(Q, M)$$
Thus, the pair $(Q, \theta^A)$
satisfies the theorem when $S = \spec A$. More generally, we can cover
$S$ by affine open subschemes. Then on their overlaps,
the resulting pairs $(\Qq,\theta)$ glue together by their uniqueness.
\end{proof}


\begin{definition}[Linear scheme and zero section]\label{nhq-linear-scheme}\source{nitsure-hilbert-quot:page-0016}
A \textbf{linear scheme} $\VV \to S$ over a noetherian base scheme $S$ is
a scheme of the form $\Spec \sym_{{\mathcal O}_S} \Qq $ where $\Qq$ is a coherent
sheaf on $S$. This is naturally a group scheme. Linear schemes
generalise the notion of (geometric) vector bundles, which are the
special case where $\Qq$ is locally free of constant rank.

The \textbf{zero section} $\VV_0 \subset \VV$ of a linear scheme
$\VV = \Spec \sym_{{\mathcal O}_S}\Qq$ is the closed subscheme defined
by the ideal generated by $\Qq$. Note that the projection $\VV_0 \to S$
is an isomorphism, and $\VV_0$ is just
the image of the zero section $0 : S \to \VV$ of the group-scheme.
\end{definition}


\begin{theorem}\label{nhq-hom-linear-scheme}\dcref{3.5}\source{nitsure-hilbert-quot:page-0016}\uses{nhq-cokernel-dual,nhq-linear-scheme} {\rm ([EGA] III 7.7.8, 7.7.9) }
Let $S$ be a noetherian scheme and $\pi : X \to S$ a projective
morphism. Let $\E$ and $\F$ be coherent sheaves on $X$.
Consider the set-valued contravariant functor $\Hom(\E,\F)$ on $S$-schemes,
which associates to any $T\to S$ the set of all ${\mathcal O}_{X_T}$-linear
homomorphisms $Hom_{X_T}(\E_T,\F_T)$ where $\E_T$ and $\F_T$ denote the
pull-backs of $\E$ and $\F$ under the projection $X_T \to X$.
If $\F$ is flat over $S$, then the above functor is representable
by a linear scheme $\VV$ over $S$.
\end{theorem}


\begin{proof} First note that if $\E$ is
a locally free ${\mathcal O}_X$-module, then $\Hom(\E,\F)$ is
the functor $T\mapsto H^0(X_T, (\F\otimes_{{\mathcal O}_X}\E^{\vee})_T)$.
The sheaf $\F\otimes_{{\mathcal O}_X}\E^{\vee}$ is again flat over $S$,
so we can apply Theorem \ref{nhq-cokernel-dual} to
get a coherent sheaf $\Qq$, such that we have
$\pi_*(\F \otimes_{{\mathcal O}_X}\E^{\vee}\otimes_{{\mathcal O}_X} \pi^*\G) =
\hom_{{\mathcal O}_S}(\Qq,\G)$ for all quasi-coherent sheaves $\G$ on $S$.
In particular, if $f: \spec R \to S$ is any morphism then
taking $\G = f_*{{\mathcal O}_R}$ we get
\begin{eqnarray*}
Mor_S(\spec R, \Spec \sym_{{\mathcal O}_S}\Qq ) & = &
Hom_{{\mathcal O}_S\text{-mod}}(\Qq,  f_*{{\mathcal O}_R}) \\
& = & H^0(X, \F \otimes_{{\mathcal O}_X}\E^{\vee}\otimes_{{\mathcal O}_X} \pi^*  f_*{{\mathcal O}_R})\\
& = & H^0(X_R, (\F \otimes_{{\mathcal O}_X}\E^{\vee})_R)\\
&= &  Hom_{X_R}(\E_R, \F_R).
\end{eqnarray*}
This shows that $\VV = \Spec \sym_{{\mathcal O}_S} \Qq $
is the required linear scheme when $\E$ is locally free on $X$.
More generally for an arbitrary coherent $\E$, over
any affine open $U\subset S$ there exist vector bundles
$E_1$ and $E_0$ on $X_U$ and a right exact sequence
$E_1 \to E_0 \to \E \to 0$. (This is where we need projectivity
of $X\to S$. Instead, we could have assumed just properness together
with the condition that locally over $S$ we have such a resolution of
$\E$.) Then applying the above argument to the functors
$\Hom(E_1,\F)$ and $\Hom(E_0,\F)$,
we get coherent sheaves $\Qq_1$ and $\Qq_0$ on
$U$, and from the natural transformation
$\Hom(E_0,\F) \to \Hom(E_1,\F)$ induced by the homomorphism
$E_1\to E_0$, we get a homomorphism $\Qq_1 \to \Qq_0$. Let
$\Qq_U$ be its cokernel, and put
$\VV_U = \Spec \Sym_{{\mathcal O}_U} \Qq_U$. It follows
from its definition (and the left exactness of Hom)
that the scheme $\VV_U$ has
the desired universal property over $U$. Therefore all such
$\VV_U$, as $U$ varies over an affine open cover of $S$, patch together
to give the desired linear scheme $\VV$.
(In sheaf terms, the sheaves $\Qq_U$ will patch together to give
a coherent sheaf $\Qq$ on $S$ with $\VV = \Spec \sym_{{\mathcal O}_S} \Qq $.)
\hfill$\square$\end{proof}





\begin{remark}\label{nhq-vanishing-scheme}\dcref{3.6}\source{nitsure-hilbert-quot:page-0017}\uses{nhq-hom-linear-scheme,nhq-linear-scheme}
In particular, note that the zero section
$\VV_0 \subset \VV$ is where
the universal homomorphism vanishes.
If $f \in Hom_{X_T}(\E_T,\F_T)$ defines a morphism
$\varphi_f : T \to \VV$, then the inverse image $f^{-1}\VV_0$
is a closed subscheme $T'$ of $T$ with the universal property
that if $U\to T$ is any morphism of schemes such that
the pull-back of $f$ is zero, then $U\to T$ factors via $T'$.






\end{remark}

\bigskip





\centerline{\large\bf Base-change for Flat Sheaves}

\medskip


The following is the main result of Grothendieck on base change for flat
families of sheaves, which is a consequence of
Theorem \ref{nhq-grothendieck-complex}.


\begin{theorem}\label{nhq-base-change-flat-sheaves}\dcref{3.7}\source{nitsure-hilbert-quot:page-0017}\uses{nhq-grothendieck-complex}
Let $\pi: X\to S$ be a proper morphism of noetherian
schemes, and let $\F$ be a coherent ${\mathcal O}_X$-module
which is flat over ${\mathcal O}_S$.
Then the following statements hold:

(1) For any integer $i$ the function $s\mapsto \dim_{\kappa(s)}H^i(X_s,\F_s)$
is upper semi-continuous on $S$,

(2) The function $s\mapsto \sum_i (-1)^i \dim_{\kappa(s)}H^i(X_s,\F_s)$
is locally constant on $S$.

(3) If for some integer $i$, there is some integer $d\ge 0$ such that
for all $s\in S$ we have $\dim_{\kappa(s)}H^i(X_s,\F_s) = d$,
then $R^i\pi_*\F$ is locally free of rank $d$,
and $(R^{i-1}\pi_*\F)_s \to H^{i-1}(X_s,\F_s)$ is an isomorphism
for all $s\in S$.

(4) If for some integer $i$ and point $s\in S$ the map
$(R^i\pi_*\F)_s \to H^i(X_s,\F_s)$ is surjective,
then there exists
an open subscheme $U\subset S$ containing $s$ such that
for any quasi-coherent ${\mathcal O}_U$-module $\G$ the natural homomorphism
$$(R^i {\pi_U}_*\F_{X_U})\otimes_{{\mathcal O}_U}\G\to
R^i{\pi_U}_*(\F_{X_U}\otimes_{{\mathcal O}_{X_U}}{\pi_U}^*\G)$$
is an isomorphism, where $X_U = \pi^{-1}(U)$ and
$\pi_U: X_U \to U$ is induced by $\pi$. In particular,
$(R^i\pi_*\F)_{s'} \to H^i(X_{s'},\F_{s'})$ is an isomorphism
for all $s'$ in $U$.


(5) If for some integer $i$ and point $s\in S$ the map
$(R^i\pi_*\F)_s \to H^i(X_s,\F_s)$ is surjective,
then the following conditions (a) and (b) are equivalent:

~~~~(a) The map $(R^{i-1}\pi_*\F)_s \to H^{i-1}(X_s,\F_s)$ is surjective.

~~~~(b) The sheaf $R^i\pi_*\F$ is locally free in a
        neighbourhood of $s$ in $S$.
\end{theorem}


\begin{proof}
\notready The source refers to Hartshorne [H], Chapter III, Section 12, for the proof of this base-change theorem; no proof is supplied here.
\end{proof}

See for example Hartshorne [H] Chapter III, Section 12 for a proof.
It is possible to replace the use of the formal function
theorem in [H] (or the original argument in [EGA] based on
completions) in proving the statement (4) above,
with an elementary argument based on applying
Nakayama lemma to the Grothendieck complex.
